
\SGARefTarget{sga3:xi:section:2}
\section*{2. Examples of Formally Smooth Functors from the Theory of
Groups of Multiplicative Type}
\addcontentsline{toc}{section}{2. Examples of Formally Smooth Functors
from the Theory of Groups of Multiplicative Type}

We shall interpret, in the language introduced in the preceding
number, the results stated in Exposé IX, no.~3, concerning
infinitesimal extensions of a homomorphism from a group of
multiplicative type.  Those results are consequences of the
vanishing of the Hochschild cohomology of such a group, established
in Exposé I.

\medskip
\noindent\SGARefTarget{sga3:xi:proposition:2:1}\textbf{Proposition 2.1.} ---
\textit{Let \(S\) be a prescheme, let \(H\) be a group of
multiplicative type over \(S\), and let \(G\) be an \(S\)-group
prescheme smooth over \(S\).  Consider the functor over \(S\)}
\SGARefTarget{sga3:xi:display:2:1:hom-functor}
\[
  M_H=\operatorname{Hom}_{S\text{-}\mathrm{gr}}(H,G).
\]
\textit{Its value at an \(S'\) over \(S\) is the set of
homomorphisms of \(S'\)-groups from \(H_{S'}\) to \(G_{S'}\).
\SGARefLink{sga3:xi:display:2:1:hom-functor}{This functor} is formally smooth over \(S\).}

Cf.\ \SGARefLink{sga3:ix:theorem:3:6}{Exposé IX, 3.6}.  More generally:

\medskip
\noindent\SGARefTarget{sga3:xi:corollary:2:2}\textbf{Corollary 2.2.} ---
\textit{Let \(S\) and \(G\) be as above, and consider a homomorphism
\(u:H_1\to H_2\) of group schemes of multiplicative type over \(S\).
With the preceding notation this induces a morphism of functors over
\(S\)}
\SGARefTarget{sga3:xi:display:2:2:restriction-map}
\[
  M_{H_2}\longrightarrow M_{H_1},
  \qquad
  M_{H_i}=\operatorname{Hom}_{S\text{-}\mathrm{gr}}(H_i,G),
\]
\textit{given by \(w\mapsto w\circ u\).  \SGARefLink{sga3:xi:display:2:2:restriction-map}{This morphism} is formally
smooth.}

Indeed, by the definitions in
\SGARefLink{sga3:xi:remark:1:6}{1.6}, this is equivalent to the
following assertion.  Suppose that \(S\) is affine and that
\(S_0\subset S\) is a subscheme defined by a nilpotent ideal.  Let
\SGARefTarget{sga3:xi:display:2:2:v}
\[
  v:H_1\longrightarrow G
\]
be a homomorphism of \(S\)-groups, and let
\SGARefTarget{sga3:xi:display:2:2:w-zero}
\[
  w_0:(H_2)_{S_0}\longrightarrow G_{S_0}
\]
be a homomorphism of \(S_0\)-groups such that
\[
  w_0u_{S_0}=v_{S_0}.
\]
Then there exists a homomorphism of \(S\)-groups
\SGARefTarget{sga3:xi:display:2:2:w}
\[
  w:H_2\longrightarrow G
\]
extending \(w_0\) and satisfying \(wu=v\).

To see this, first extend \(w_0\) to a homomorphism of \(S\)-groups
\[
  w':H_2\longrightarrow G,
\]
which is possible by
\SGARefLink{sga3:xi:proposition:2:1}{2.1}.  Set
\[
  v'=w'u:H_1\longrightarrow G.
\]
By the hypothesis on \(w_0\), one has \(v'_{S_0}=v_{S_0}\).
Consequently, by Exposé IX, 3.6, there exists an element
\(g\in G(S)\) whose image in \(G(S_0)\) is the identity element and
such that
\[
  v=\operatorname{int}(g)v'.
\]
It follows that
\[
  v=\bigl(\operatorname{int}(g)w'\bigr)u,
\]
and it is therefore enough to take
\[
  w=\operatorname{int}(g)w'.
\]

\medskip
\noindent\SGARefTarget{sga3:xi:corollary:2:3}\textbf{Corollary 2.3.} ---
\textit{With the notation of
\SGARefLink{sga3:xi:proposition:2:1}{2.1}, regard \(M_H=M\) as a
functor with operator group \(G\), where \(G\) acts by
\[
  (v,g)\longmapsto\operatorname{int}(g)\circ v.
\]
Then the corresponding morphism}
\SGARefTarget{sga3:xi:display:2:3:action-map}
\[
  R:G\times_S M\longrightarrow M\times_S M,
\]
\textit{defined by
\[
  (g,v)\longmapsto\bigl(\operatorname{int}(g)\circ v,v\bigr),
\]
is formally smooth.}

After a base change \(S'\to S\), this is equivalent to the following
assertion.

\medskip
\noindent\SGARefTarget{sga3:xi:corollary:2:4}\textbf{Corollary 2.4.} ---
\textit{With the notation of
\SGARefLink{sga3:xi:proposition:2:1}{2.1}, let
\(v_1,v_2:H\to G\) be two morphisms of \(S\)-groups, and let
\(\operatorname{Transp}(v_1,v_2)\) be the subfunctor of \(G\)
consisting of the \(g\) such that
\[
  \operatorname{int}(g)\circ v_1=v_2.
\]
Then this functor is formally smooth over \(S\).  In particular, if
\(v_1=v_2=v\), the functor \(\operatorname{Centr}(v)\), the subgroup
of \(G\) consisting of the \(h\) such that
\(\operatorname{int}(h)v=v\), is formally smooth over \(S\).}

\emph{N.B.} The pair \((v_1,v_2)\) may be regarded as a section over
\(S\) of the second factor in the codomain of the morphism \(R\) in
\SGARefLink{sga3:xi:corollary:2:3}{2.3}, and
\(\operatorname{Transp}(v_1,v_2)\) as the inverse-image functor of
that section under \(R\).

The assertion in
\SGARefLink{sga3:xi:corollary:2:4}{2.4} is itself equivalent to the
following.  Suppose that \(S\) is affine and that \(S_0\subset S\) is
a subscheme defined by a nilpotent ideal.  If
\(g_0\in G(S_0)\) satisfies
\[
  \operatorname{int}(g_0)(v_1)_{S_0}=(v_2)_{S_0},
\]
then \(g_0\) extends to an element \(g\in G(S)\) such that
\[
  \operatorname{int}(g)v_1=v_2.
\]
To prove this, first extend \(g_0\) to a section \(g'\) of \(G\) over
\(S\), which is possible because \(G\) is smooth over \(S\), and put
\[
  v'_2=\operatorname{int}(g')v_1.
\]
The morphisms \(v_2\) and \(v'_2\) have the same restriction over
\(S_0\).  Hence, by \SGARefLink{sga3:ix:theorem:3:6}{Exposé IX, 3.6}, there exists
\(g''\in G(S)\) inducing the identity section over \(S_0\) and such
that
\[
  v_2=\operatorname{int}(g'')v'_2.
\]
Therefore
\[
  v_2
  =\operatorname{int}(g'')\operatorname{int}(g')v_1
  =\operatorname{int}(g''g')v_1,
\]
and it is enough to take \(g=g''g'\).

\medskip
\noindent\SGARefTarget{sga3:xi:proposition:2:1:bis}\textbf{Proposition 2.1 bis.} ---
\textit{Let \(S\) be a prescheme and let \(G\) be an \(S\)-group
prescheme smooth over \(S\).  Consider the functor}
\SGARefTarget{sga3:xi:display:2:1:bis:functor}
\[
  M:(\mathbf{Sch})^\circ/S\longrightarrow(\mathbf{Ens}),
\]
\textit{defined by}
\SGARefTarget{sga3:xi:display:2:1:bis:value}
\[
  M(S')
  =\left\{
    \begin{array}{c}
      \text{multiplicative-type subgroups}\\[-2pt]
      \text{of \(G_{S'}\)}
    \end{array}
  \right\}.
\]
\textit{Then \(M\) is formally smooth over \(S\).}

Cf.\ \SGARefLink{sga3:ix:theorem:3:6:bis}{Exposé IX, 3.6 bis}.

\medskip
\noindent\SGARefTarget{sga3:xi:corollary:2:2:bis}\textbf{Corollary 2.2 bis.} ---
\textit{Let \(n\) be an integer, and consider the morphism of
functors}
\SGARefTarget{sga3:xi:display:2:2:bis:phi}
\[
  \varphi_n:M\longrightarrow M
\]
\textit{defined by}
\SGARefTarget{sga3:xi:display:2:2:bis:kernel}
\[
  \varphi_n(H)={}_nH=\ker\!\bigl(n\cdot\operatorname{id}_H\bigr).
\]
\textit{Then \(\varphi_n\) is formally smooth.  For every integer
\(p>0\), let \(M_p\) denote the subfunctor of \(M\) for which
\(M_p(S')\) is the set of multiplicative-type subgroups \(H\) of
\(G_{S'}\) satisfying \({}_pH=H\).  Then the morphism induced by
\(\varphi_n\)}
\SGARefTarget{sga3:xi:display:2:2:bis:restricted-map}
\[
  M_{np}\longrightarrow M_n
\]
\textit{is formally smooth.}

The second assertion is trivially contained in the first and is
included only for the convenience of a later reference.  The proof
of the first assertion is analogous to that of
\SGARefLink{sga3:xi:corollary:2:2}{2.2}, now invoking \SGARefLink{sga3:ix:theorem:3:6:bis}{Exposé IX,
3.6 bis}.

\medskip
\noindent\SGARefTarget{sga3:xi:corollary:2:3:bis}\textbf{Corollary 2.3 bis.} ---
\textit{With the notation of
\SGARefLink{sga3:xi:proposition:2:1:bis}{2.1 bis}, regard \(M\) as a
functor with operator group \(G\), where \(G\) acts by
\[
  (g,H)\longmapsto\operatorname{int}(g)(H).
\]
Then the corresponding morphism}
\SGARefTarget{sga3:xi:display:2:3:bis:action-map}
\[
  R:G\times_S M\longrightarrow M\times_S M,
\]
\textit{defined by
\[
  (g,v)\longmapsto\bigl(\operatorname{int}(g)v,v\bigr),
\]
is formally smooth.}

This is equivalent to the following assertion.

\medskip
\noindent\SGARefTarget{sga3:xi:corollary:2:4:bis}\textbf{Corollary 2.4 bis.} ---
\textit{With the notation of
\SGARefLink{sga3:xi:proposition:2:1:bis}{2.1 bis}, let \(H_1\) and
\(H_2\) be two multiplicative-type subgroups of \(G\), and let
\(\operatorname{Transp}_G(H_1,H_2)\) be the subfunctor of \(G\)
consisting of the \(g\) such that
\[
  \operatorname{int}(g)H_1=H_2.
\]
Then this functor is formally smooth over \(S\).  In particular, if
\(H_1=H_2=H\), the subfunctor
\(\operatorname{Norm}_G(H)\) of \(G\), the normalizer of \(H\), is
formally smooth over \(S\).}

The proof is analogous to that of
\SGARefLink{sga3:xi:corollary:2:4}{2.4}, again invoking \SGARefLink{sga3:ix:theorem:3:6:bis}{Exposé IX,
3.6 bis}.

\medskip
\noindent\SGARefTarget{sga3:xi:proposition:2:5}\textbf{Proposition 2.5.} ---
\textit{Let \(S\) be a prescheme, let \(G\) be an \(S\)-group
prescheme smooth over \(S\), let \(K\) be an \(S\)-group
subprescheme that is smooth over \(S\) or of multiplicative type,
let \(H\) be an \(S\)-group prescheme of multiplicative type, and let
\(u:H\to G\) be a homomorphism of \(S\)-groups.  Denote by
\(\operatorname{Transp}_G(u,K)\) the subfunctor of \(G\) whose value
at an \(S'\) over \(S\) consists of the elements \(g\in G(S')\) for
which}
\SGARefTarget{sga3:xi:display:2:5:transporter-condition}
\[
  \operatorname{int}(g)u_{S'}:
  H_{S'}\longrightarrow G_{S'}
\]
\textit{factors through \(K_{S'}\).  Then this functor is formally
smooth over \(S\).}

The proof is analogous to that of
\SGARefLink{sga3:xi:corollary:2:2}{2.2}, invoking \SGARefLink{sga3:ix:theorem:3:6}{Exposé IX, 3.6} and
\SGARefLink{sga3:x:proposition:2:1}{Exposé X, 2.1}, the latter in the case where \(K\) is of
multiplicative type.

